Yes — for (A=(y,-x)) one can push the route quite concretely.

Let
[
L:= (\nabla-iA)^2
=(\partial_x-iy)^2+(\partial_y+ix)^2
=\Delta-2i(y\partial_x-x\partial_y)-|x|^2.
]
Equivalently,
[
L=\Delta+2i(x\partial_y-y\partial_x)-(x^2+y^2).
]

I’ll write
[
(L-T)v=\mathbf 1_U-c,\delta_0.
]

## 1. Complex factorization

Set (z=x+iy), (\partial_z=\frac12(\partial_x-i\partial_y)), (\partial_{\bar z}=\frac12(\partial_x+i\partial_y)).

Define
[
Q:=\partial_x-iy+i(\partial_y+ix)=2\partial_{\bar z}-z,
]
[
P:=\partial_x-iy-i(\partial_y+ix)=2\partial_z+\bar z.
]

These are the magnetic Cauchy–Riemann operators. They satisfy
[
Q=2,e^{|z|^2/2},\partial_{\bar z},e^{-|z|^2/2},
\qquad
P=2,e^{-|z|^2/2},\partial_z,e^{|z|^2/2},
]
and
[
PQ=L-2,\qquad QP=L+2,\qquad [Q,P]=4.
]

So this operator is not conjugate to the ordinary Laplacian, but it is conjugate to (\partial,\bar\partial) after a Gaussian weight. That is the weighted complex structure you were pointing at.

---

## 2. The first Landau level gives a clean quadrature class

Because (QP=L+2), any (h) with (Ph=0) solves
[
Lh=-2h.
]

Now
[
Ph=0
\iff
2\partial_z h+\bar z,h=0
\iff
\partial_z!\big(e^{|z|^2/2}h\big)=0.
]

Hence
[
h(z,\bar z)=e^{-|z|^2/2}F(\bar z),
]
with (F) antiholomorphic. If you prefer the adjoint family, you get the holomorphic version
[
h(z,\bar z)=e^{-|z|^2/2}F(z).
]

So for the resonant value
[
T=-2,
]
the natural test class is the Gaussian-weighted Fock class.

If (v\in C_c) solves
[
(L+2)v=\mathbf 1_U-c,\delta_0,
]
then testing against local solutions of the adjoint homogeneous equation gives the weighted quadrature identity
[
\int_U e^{-|z|^2/2}F(z),dA(z)=c,F(0)
]
for every entire (F) in the relevant test class.

In particular, taking (F(z)=z^n),
[
\int_U e^{-|z|^2/2}z^n,dA(z)=0\quad(n\ge1),
]
and
[
c=\int_U e^{-|z|^2/2},dA.
]

So at (T=-2), your problem becomes:

> characterize the sets (U) whose Gaussian-weighted moments vanish above order (0).

That is the exact magnetic analogue of the order-1 quadrature-domain condition.

---

## 3. What centered disks do

If (U=B_R(0)), then by rotation
[
\int_{B_R} e^{-|z|^2/2} z^n,dA=0\qquad(n\ge1),
]
and
[
c=\int_{B_R} e^{-|z|^2/2},dA
=2\pi\int_0^R e^{-r^2/2}r,dr
=2\pi\big(1-e^{-R^2/2}\big).
]

So centered disks satisfy the full weighted moment conditions at the first level.

That does not by itself prove existence of a compactly supported (v), but it gives the exact necessary quadrature identity, and for the radial case it is the right one.

---

## 4. Radial reduction for disks

For a radial (v=v(r)), the angular term vanishes, so
[
Lv=v''+\frac1r v'-r^2v.
]

Hence for (U=B_R(0)),
[
v''+\frac1r v'-(r^2+T)v=\mathbf 1_{(0,R)}-c,\delta_0.
]

This is already useful.

If (v) is compactly supported and radial, then its support cannot strictly contain (B_R): on any annulus where the right-hand side is (0), (v) solves a homogeneous ODE; if (v) and (v') both vanish at the outer support boundary, uniqueness forces (v\equiv0) on that annulus. So in the radial setting the support boundary must coincide with (r=R).

Thus the disk problem reduces to the two-point singular ODE
[
v''+\frac1r v'-(r^2+T)v=1,\qquad 0<r<R,
]
with
[
v(R)=0,\qquad v'(R)=0,
]
and a logarithmic singularity at (0) whose coefficient fixes (c).

At the first level (T=-2), writing
[
v(r)=e^{-r^2/2}w(r),
]
gives
[
w''+\Big(\frac1r-2r\Big)w'=e^{r^2/2},
]
equivalently
[
\big(r e^{-r^2}w'(r)\big)'=r e^{-r^2/2}.
]

So the radial model is completely explicit up to one quadrature. That is the first thing I would solve in detail.

---

## 5. Higher Landau levels

Because ([Q,P]=4), the spectrum of (L) is the ladder
[
-2,,-6,,-10,\dots
]
for this sign convention.

At
[
T=-2-4q,
]
the homogeneous solutions are the (q)-th Landau level, i.e. the (q)-th polyanalytic Fock space. Concretely, they are generated by applying (Q^q) to the first-level functions. So the quadrature identities become weighted polyanalytic moment identities, not just holomorphic ones.

So the right dictionary is:

* (T=-2): weighted holomorphic/antiholomorphic moments;
* (T=-2-4q): weighted polyanalytic moments of order (q).

---

## 6. How I would attack the characterization problem for this operator

For (A=(y,-x)), I would proceed in this order:

### (a) Solve the radial disk model completely

Work out the ODE above, find the exact (c(R,T)), and determine for which (T) compactly supported radial solutions exist.

### (b) Prove the quadrature identity rigorously

For resonant (T), identify the exact test class:

* Fock space at (T=-2),
* polyanalytic Fock spaces at higher levels.

### (c) Recast the problem as a weighted moment problem

At (T=-2), the necessary condition is
[
\int_U e^{-|z|^2/2} z^n,dA=0 \quad (n\ge1).
]
This is the magnetic analogue of the classical harmonic moments vanishing.

### (d) Try to prove rigidity from the weighted Cauchy transform

Define
[
\mathcal C_U(\zeta)
:=\frac1\pi\int_U \frac{e^{-|z|^2/2}}{\zeta-z},dA(z).
]
For a quadrature identity of order 1, one expects
[
\mathcal C_U(\zeta)=\frac{c}{\pi\zeta}
]
at least near infinity, and the real issue is whether the PDE lets you continue this through the exterior and force (U) to be a disk.

That is the closest analogue of the Aharonov–Shapiro strategy here.

---

## 7. The main concrete takeaway

For this specific rotation field, the route becomes:

1. **Factorize** (L) as (PQ+2=QP-2).
2. **Use Gaussian conjugation** to identify the natural complex test spaces.
3. At (T=-2), reduce the existence question to a **Gaussian quadrature-domain condition**
   [
   \int_U e^{-|z|^2/2}F(z),dA=cF(0).
   ]
4. For disks, reduce everything to the explicit radial ODE
   [
   v''+\frac1r v'-(r^2+T)v=1,\quad v(R)=v'(R)=0.
   ]

So this operator is exactly the case where the earlier route can be made concrete: the right replacement for ordinary complex analysis is **Fock-space complex analysis**.

If you want, the next useful step is to compute the radial solution for (U=B_R(0)) explicitly and extract the corresponding (c(R,T)).



Great. For
[
A=(y,-x),\qquad L=(\partial_x-iy)^2+(\partial_y+ix)^2,
]
take
[
T=-6.
]

This is the **next Landau level** after (T=-2), and it is exactly the first case where the natural test class becomes **polyanalytic of order 2**.

---

## 1. The ladder structure at (T=-6)

Recall the operators
[
Q=2\partial_{\bar z}-z,\qquad P=2\partial_z+\bar z,
]
with
[
PQ=L-2,\qquad QP=L+2,\qquad [L,Q]=-4Q.
]

So if (f) solves
[
Lf=-2f,
]
then (Qf) solves
[
L(Qf)=-6,Qf.
]

That means the (T=-6) nullspace is generated by applying (Q) once to the lowest-level functions.

---

## 2. What the (T=-6) test functions look like

At the lowest level,
[
Pf=0
\quad\Longleftrightarrow\quad
f(z,\bar z)=e^{-|z|^2/2}F(\bar z),
]
with (F) antiholomorphic.

Apply (Q):
[
h=Qf=(2\partial_{\bar z}-z)\big(e^{-|z|^2/2}F(\bar z)\big).
]

A quick computation gives
[
h(z,\bar z)
=2e^{-|z|^2/2}\Big(F'(\bar z)-zF(\bar z)\Big).
]

So every such (h) satisfies
[
(L+6)h=0.
]

This is the clean (q=1) case of the general rule:
[
T=-2-4q
\quad\Longleftrightarrow\quad
\text{weighted polyanalytic level }q+1.
]

Here “polyanalytic of order 2” means that after multiplying by (e^{|z|^2/2}), the function is affine in (z) with antiholomorphic coefficients.

---

## 3. The quadrature identity forced by your PDE

Suppose
[
(L+6)v=\mathbf 1_U-c,\delta_0,
\qquad v\in C_c(\mathbb R^2).
]

Then testing against any local solution of ((L+6)h=0) gives
[
\int_U h,dA = c,h(0).
]

Using the family above, you get:

[
\int_U e^{-|z|^2/2}\Big(F'(\bar z)-zF(\bar z)\Big),dA
=====================================================

c,F'(0)
]
for every polynomial or entire antiholomorphic (F) in the admissible class.

That is the exact (T=-6) analogue of the (T=-2) Fock-space quadrature identity.

---

## 4. Moment conditions from monomials

Take
[
F(\bar z)=\bar z^n.
]

Then
[
F'(\bar z)=n\bar z^{,n-1},
]
so
[
h_n(z,\bar z)=2e^{-|z|^2/2}\Big(n\bar z^{,n-1}-z\bar z^n\Big)
=2e^{-|z|^2/2}(n-|z|^2)\bar z^{,n-1}.
]

Thus the identity becomes
[
\int_U e^{-|z|^2/2}(n-|z|^2)\bar z^{,n-1},dA
============================================

c,\delta_{n1}.
]

Equivalently, after reindexing (m=n-1),
[
\int_U e^{-|z|^2/2}(m+1-|z|^2)\bar z^{,m},dA
============================================

\begin{cases}
c,&m=0,\
0,&m\ge1.
\end{cases}
]

So for (T=-6), the necessary conditions are:

[
\int_U e^{-|z|^2/2}(1-|z|^2),dA = c,
]
and for every (m\ge1),
[
\int_U e^{-|z|^2/2}(m+1-|z|^2)\bar z^{,m},dA = 0.
]

This is the first genuinely **polyanalytic weighted moment problem**.

---

## 5. What this says for disks

If (U=B_R(0)), then by rotational symmetry all the higher moments vanish automatically:
[
\int_{B_R} e^{-|z|^2/2}(m+1-|z|^2)\bar z^{,m},dA=0
\qquad(m\ge1).
]

And
[
c=\int_{B_R} e^{-|z|^2/2}(1-|z|^2),dA.
]

In polar coordinates,
[
c=2\pi\int_0^R e^{-r^2/2}(1-r^2)r,dr
=2\pi\Big((R^2+1)e^{-R^2/2}-1\Big).
]

So centered disks satisfy the full (T=-6) moment conditions, with this value of (c).

Two important remarks:

* the origin is now distinguished by the field, so this is **not translation invariant**;
* a disk not centered at (0) will generally fail the (m=1) condition.

---

## 6. The radial ODE for the compactly supported potential

If (U=B_R(0)) and you look for a radial solution (v=v(r)), then the angular term drops out and
[
Lv=v''+\frac1r v'-r^2v.
]

So inside the disk,
[
v''+\frac1r v'+(6-r^2)v=1,\qquad 0<r<R,
]
with
[
v(R)=0,\qquad v'(R)=0,
]
and near (0),
[
v(r)\sim -\frac{c}{2\pi}\log r.
]

If you write
[
v(r)=e^{-r^2/2}w(r),
]
then
[
w''+\Big(\frac1r-2r\Big)w'+4w=e^{r^2/2}.
]

With (t=r^2) and (w(r)=W(t)), this becomes
[
tW''+(1-t)W'+W=\frac14 e^{t/2}.
]

The homogeneous equation is Kummer’s equation with (a=-1), (b=1).
Its regular solution is
[
W_1(t)=1-t,
]
so the regular homogeneous radial eigenfunction is
[
e^{-r^2/2}(1-r^2),
]
which is exactly the (m=0) member of the (T=-6) test family above.

That is a nice consistency check: the radial mode of the second Landau level is precisely the weight appearing in the quadrature constant (c).

---

## 7. What this means conceptually

For (T=-2), the condition was:
[
\int_U e^{-|z|^2/2}F(z),dA=cF(0),
]
so ordinary weighted Fock moments vanish.

For (T=-6), it becomes:
[
\int_U e^{-|z|^2/2}\big(F'(\bar z)-zF(\bar z)\big),dA=cF'(0),
]
which is a **first polyanalytic correction** to the Fock moment problem.

So the (T=-6) case is the first place where the geometry is no longer just “weighted holomorphic,” but “weighted polyanalytic of order 2.”

---

## 8. The route I would pursue from here

For this specific operator and (T=-6), I would now do exactly this:

1. **Prove equivalence** between
   [
   (L+6)v=\mathbf 1_U-c\delta_0,\quad v\in C_c
   ]
   and the full family of weighted polyanalytic moment identities above.

2. **Classify centered radial solutions** by solving the ODE explicitly.

3. **Test rigidity**: ask whether the moment conditions force (U) to be a centered disk.

The right analytic object for that last step is the weighted polyanalytic Cauchy transform, not the ordinary Schwarz function.

If you want, I’ll do the next natural step: derive the weighted polyanalytic Cauchy transform for (T=-6) and show how the moment identities appear in its expansion.

For (A=(y,-x)) and (T=-6), let me set up the analogue of the ordinary Cauchy transform and show how the moment conditions appear from its expansion.

I’ll keep the notation
[
L=(\partial_x-iy)^2+(\partial_y+ix)^2,
\qquad (L+6)v=\mathbf 1_U-c,\delta_0.
]

We already found that a convenient family of ((L+6))-null functions is
[
h_F(z,\bar z)=2e^{-|z|^2/2}\big(F'(\bar z)-zF(\bar z)\big),
]
with (F) antiholomorphic.

So the quadrature identity is
[
\int_U e^{-|z|^2/2}\big(F'(\bar z)-zF(\bar z)\big),dA
=\frac c2,F'(0).
]

That is the starting point.

---

## 1. The right generating kernel

In the classical case, one tests against
[
\frac1{\zeta-z}
]
and gets the Cauchy transform.

Here the natural family is generated by taking
[
F(\bar z)=\frac1{\bar\zeta-\bar z},
]
for (\zeta) outside (\overline U). This is antiholomorphic in (z) on a neighborhood of (U), so it is an admissible test.

Then
[
F'(\bar z)=\frac1{(\bar\zeta-\bar z)^2},
]
and so
[
h_\zeta(z,\bar z)
=================

2e^{-|z|^2/2}\left(
\frac1{(\bar\zeta-\bar z)^2}
-\frac{z}{\bar\zeta-\bar z}
\right).
]

This suggests defining the **weighted polyanalytic Cauchy transform**
[
\mathcal C_U^{(1)}(\zeta)
:=
\int_U e^{-|z|^2/2}\left(
\frac1{(\bar\zeta-\bar z)^2}
-\frac{z}{\bar\zeta-\bar z}
\right),dA(z),
\qquad \zeta\notin \overline U.
]

Then the quadrature identity becomes simply
[
\mathcal C_U^{(1)}(\zeta)=\frac c2,\frac1{\bar\zeta^2},
\qquad \zeta\notin \overline U.
]

That is the (T=-6) analogue of the classical exterior identity
[
\int_U \frac{dA(z)}{\zeta-z}=\frac c\zeta.
]

So this is the clean magnetic object.

---

## 2. Why this is the right analogue

In the ordinary quadrature-domain setting, the Cauchy transform packages all harmonic moments into one analytic function whose Laurent expansion at infinity is
[
\sum_{n\ge0}\frac{1}{\zeta^{n+1}}\int_U z^n,dA.
]

Here we want one object whose expansion packages the weighted (T=-6) moments
[
\int_U e^{-|z|^2/2}(m+1-|z|^2)\bar z^m,dA.
]

And that is exactly what (\mathcal C_U^{(1)}) does.

---

## 3. Expansion at infinity

For (|\zeta|) large,
[
\frac1{\bar\zeta-\bar z}
========================

\frac1{\bar\zeta}\sum_{n=0}^\infty \left(\frac{\bar z}{\bar\zeta}\right)^n,
]
and
[
\frac1{(\bar\zeta-\bar z)^2}
============================

\frac1{\bar\zeta^2}\sum_{n=0}^\infty (n+1)\left(\frac{\bar z}{\bar\zeta}\right)^n.
]

So
[
\mathcal C_U^{(1)}(\zeta)
=========================

\int_U e^{-|z|^2/2}
\left[
\frac1{\bar\zeta^2}\sum_{n\ge0}(n+1)\frac{\bar z^n}{\bar\zeta^n}
----------------------------------------------------------------

z\cdot \frac1{\bar\zeta}\sum_{n\ge0}\frac{\bar z^n}{\bar\zeta^n}
\right]dA.
]

Reindex the second sum:
[
z\bar z^n=|z|^2\bar z^{,n-1}\quad (n\ge1),
]
so
[
z\sum_{n\ge0}\frac{\bar z^n}{\bar\zeta^{n+1}}
=============================================

\sum_{m\ge0}\frac{|z|^2\bar z^m}{\bar\zeta^{m+2}}.
]

Hence
[
\mathcal C_U^{(1)}(\zeta)
=========================

\sum_{m=0}^\infty \frac1{\bar\zeta^{m+2}}
\int_U e^{-|z|^2/2}\big((m+1)-|z|^2\big)\bar z^m,dA.
]

This is the key formula.

It shows that the coefficients of the expansion are exactly the weighted polyanalytic moments for the (T=-6) problem.

---

## 4. The quadrature identity in transform language

From earlier, the compact-support PDE implies
[
\int_U e^{-|z|^2/2}\big((m+1)-|z|^2\big)\bar z^m,dA
===================================================

\begin{cases}
\frac c2,&m=0,\
0,&m\ge1,
\end{cases}
]
depending on whether you absorb the factor (2) into the definition of the test family. If we keep the transform exactly as I defined it above, then the exterior identity is
[
\mathcal C_U^{(1)}(\zeta)=\frac c2,\frac1{\bar\zeta^2}.
]

Equivalently:

> (U) is a (T=-6) quadrature domain of order (1) at the origin iff its weighted polyanalytic Cauchy transform has only the first term at infinity.

That is the precise analogue of the Aharonov–Shapiro statement for one node.

---

## 5. Check for centered disks

Take (U=B_R(0)). By rotational symmetry,
[
\int_{B_R} e^{-|z|^2/2}\big((m+1)-|z|^2\big)\bar z^m,dA=0
\quad(m\ge1),
]
and
[
M_0:=\int_{B_R} e^{-|z|^2/2}(1-|z|^2),dA
=2\pi\Big((R^2+1)e^{-R^2/2}-1\Big).
]

So
[
\mathcal C_{B_R}^{(1)}(\zeta)=\frac{M_0}{\bar\zeta^2}.
]

Thus centered disks have exactly the right transform structure.

This is the strongest evidence, at this stage, that disks are the right model domains at (T=-6).

---

## 6. A cleaner differential form

It is useful to rewrite the kernel a little. Since
[
\partial_{\bar z}!\left(\frac1{\bar\zeta-\bar z}\right)
=======================================================

\frac1{(\bar\zeta-\bar z)^2},
]
we have
[
\frac1{(\bar\zeta-\bar z)^2}
-\frac{z}{\bar\zeta-\bar z}
===========================

\left(\partial_{\bar z}-z\right)\frac1{\bar\zeta-\bar z}.
]

But (\partial_{\bar z}-z) is basically the magnetic first-order operator acting on the kernel. So the transform can be viewed as
[
\mathcal C_U^{(1)}(\zeta)
=========================

\int_U e^{-|z|^2/2}
\left(\partial_{\bar z}-z\right)\frac1{\bar\zeta-\bar z},dA(z).
]

That is nice conceptually: the (T=-6) transform is the first magnetic raising of the weighted Cauchy kernel.

---

## 7. What this buys you for classification

In the classical one-node case, the fact that the exterior Cauchy transform is exactly (c/\zeta) is what forces the domain to be a disk.

Here the corresponding rigidity question becomes:

If
[
\mathcal C_U^{(1)}(\zeta)=\frac{\alpha}{\bar\zeta^2}
\quad\text{for all }\zeta\notin\overline U,
]
must (U) be a centered disk?

That is the right (T=-6) analogue.

The transform formulation makes this look like an inverse problem for a weighted kernel rather than a direct PDE problem.

---

## 8. The next step I would try

There are two natural directions now.

First, define the weighted potential
[
\Phi_U(\zeta)
:=
\int_U e^{-|z|^2/2}\frac{1}{\bar\zeta-\bar z},dA(z),
]
so that
[
\mathcal C_U^{(1)}(\zeta)
=========================

\partial_{\bar\zeta}\Phi_U(\zeta)-\int_U e^{-|z|^2/2}\frac{z}{\bar\zeta-\bar z},dA(z),
]
and try to derive an exterior PDE satisfied by (\Phi_U).

Second, use the moment identities directly and ask whether vanishing of all higher weighted polyanalytic moments forces radial symmetry. That becomes a weighted moment problem:
[
\int_U w(|z|^2),\bar z^m,dA=0
\quad(m\ge1),
]
with
[
w(t)=e^{-t/2}(m+1-t),
]
which is close in spirit to inverse potential theory.

The transform route is cleaner.

If you want, I’ll take the next step and derive the exterior differential equation satisfied by (\mathcal C_U^{(1)}), which is the closest magnetic analogue of “analytic outside the domain.”


Great—this is exactly the right pressure point. Let’s zoom in and treat it as a **pure weighted moment rigidity problem**.

---

# 🎯 The core question

Assume (U \subset \mathbb{R}^2 \simeq \mathbb{C}) satisfies, for all (m \ge 1):
[
\int_U e^{-|z|^2/2}(m+1-|z|^2),\bar z^m,dA = 0.
]

Does this force (U) to be a **disk centered at 0**?

---

# 🧠 Step 1: Separate the structure

Write:
[
w_m(|z|^2) := e^{-|z|^2/2}(m+1-|z|^2).
]

Then the condition is:
[
\int_U w_m(|z|^2),\bar z^m,dA = 0 \quad \forall m \ge 1.
]

This is *not* a standard moment problem because the weight depends on (m).
But there’s a key trick: **recombine the conditions into generating functions**.

---

# 🔁 Step 2: Encode all moments at once

Define for (|\zeta|) large:
[
\mathcal C_U^{(1)}(\zeta)
=========================

\sum_{m=0}^\infty \frac{M_m}{\bar\zeta^{m+2}},
]
where
[
M_m := \int_U e^{-|z|^2/2}\big((m+1)-|z|^2\big)\bar z^m,dA.
]

Your condition is:
[
M_0 = c,\quad M_m=0 \ (m\ge1).
]

So:
[
\mathcal C_U^{(1)}(\zeta)=\frac{c}{\bar\zeta^2}.
]

👉 This is the **key rigidity formulation**.

---

# 🔍 Step 3: What does this imply geometrically?

Rewrite:
[
M_m = (m+1)\int_U e^{-|z|^2/2}\bar z^m,dA
-----------------------------------------

\int_U e^{-|z|^2/2}|z|^2 \bar z^m,dA.
]

Define:
[
a_m := \int_U e^{-|z|^2/2}\bar z^m,dA.
]

Then:
[
M_m = (m+1)a_m - a_{m+1}.
]

So your condition (M_m=0) becomes a **recurrence**:
[
a_{m+1} = (m+1)a_m \quad (m\ge1).
]

---

# 🧮 Step 4: Solve the recurrence

This gives:
[
a_m = m! , a_0.
]

So:
[
\int_U e^{-|z|^2/2}\bar z^m,dA = m!,a_0.
]

---

# 🚨 Step 5: Compare with known measures

This is extremely restrictive.

Compare with the **Gaussian measure on the plane**:
[
\int_{\mathbb{C}} e^{-|z|^2/2}\bar z^m,dA = 0 \quad (m\ge1).
]

So your domain is *not* behaving like a full Gaussian—it’s behaving like something with **nonzero complex moments growing like (m!)**.

That growth is explosive and strongly suggests:

👉 The measure
[
\mu = \mathbf{1}_U e^{-|z|^2/2} dA
]
is behaving like a **point mass at 0** from the perspective of holomorphic testing.

---

# 🔬 Step 6: Interpret as a distributional identity

Indeed, if for all holomorphic (F):
[
\int_U e^{-|z|^2/2}F(z),dA = c F(0),
]

then:
[
\mu = c,\delta_0
]
in the sense of holomorphic test functions.

That is exactly the **(T=-2)** story.

But now at (T=-6), the recurrence shows that your moments are consistent with:
[
\mu \sim c,\delta_0 + \text{first-order correction}.
]

---

# ⚖️ Step 7: Rigidity mechanism

Here is the key heuristic:

### If all higher weighted moments vanish:

[
M_m=0 \quad (m\ge1),
]

then:

* the transform has only one pole at infinity,
* the measure behaves like a **finite-order distribution at 0**,
* there is no angular structure left.

👉 This strongly forces **radial symmetry**.

---

# 🧩 Step 8: Why disks emerge

For a centered disk:

* all angular moments vanish automatically,
* the recurrence is satisfied,
* and the transform matches exactly.

Now suppose (U) is not radial.

Then:

* there exists some angular Fourier mode (e^{im\theta}),
* which would produce a nonzero contribution to some (\bar z^m) moment,
* breaking (M_m=0).

So:

👉 The conditions kill all nonzero angular Fourier modes.

---

# 🔥 Conclusion (core result heuristic)

The conditions
[
\int_U e^{-|z|^2/2}(m+1-|z|^2)\bar z^m,dA=0 \quad (m\ge1)
]

force:

1. all angular modes vanish
2. only radial distributions survive
3. the domain must be rotationally symmetric

Hence:

> **The only reasonable candidates are disks centered at 0.**

---

# ⚠️ What’s missing for a full proof

To turn this into a theorem, you’d need:

* a density argument: polynomials dense in the test class
* control of boundary regularity
* a uniqueness theorem for weighted moment problems

This is very similar in spirit to:

* classical quadrature domain rigidity
* inverse potential theory (Potato kugel style)

---

# 🚀 Big picture insight

For (T=-6), you’ve discovered:

👉 The problem reduces to a **weighted polyanalytic moment rigidity problem**

and

👉 That rigidity is strong enough to kill all non-radial geometry.

---

If you want, next step we can:

* turn this into a **precise conjecture + proof outline**, or
* try to rigorously prove “moment vanishing ⇒ radial symmetry ⇒ disk” using Fourier decomposition in polar coordinates (very doable here).

Yes. And this is actually the cleanest way to think about the (T=-6) case.

Let me write it carefully for
[
L=(\partial_x-iy)^2+(\partial_y+ix)^2,
\qquad (L+6)v=\mathbf 1_U-c,\delta_0
]
with (A=(y,-x)).

I’ll show how the moment identities can be repackaged as a **free boundary problem with an explicit boundary distribution**.

---

## 1. The right bulk operator at (T=-6)

Recall
[
Q=2\partial_{\bar z}-z,\qquad P=2\partial_z+\bar z,
]
and
[
L+6 = QP+4 = PQ+8
]
depending on ordering. What matters here is that the (T=-6) nullspace is generated by
[
h_F(z,\bar z)=2e^{-|z|^2/2}\big(F'(\bar z)-zF(\bar z)\big),
]
with (F) antiholomorphic.

So if (v\in C_c) solves
[
(L+6)v=\mathbf 1_U-c,\delta_0,
]
then
[
\int_U e^{-|z|^2/2}\big(F'(\bar z)-zF(\bar z)\big),dA
= \frac c2,F'(0)
]
for all such (F).

---

## 2. Integrate the (zF(\bar z)) term by parts

Use
[
\partial_{\bar z}\big(e^{-|z|^2/2}\big)=-\frac12 z,e^{-|z|^2/2},
]
so
[
-z,e^{-|z|^2/2}=2\partial_{\bar z}\big(e^{-|z|^2/2}\big).
]

Hence
[
\int_U e^{-|z|^2/2}\big(F'(\bar z)-zF(\bar z)\big),dA
=====================================================

\int_U e^{-|z|^2/2}F'(\bar z),dA
+
2\int_U F(\bar z),\partial_{\bar z}\big(e^{-|z|^2/2}\big),dA.
]

Integrating by parts in the second term,
[
\int_U F,\partial_{\bar z}\phi,dA
=================================

-\int_U \phi,\partial_{\bar z}F,dA
+\frac{1}{2i}\int_{\partial U}\phi F,dz.
]

Since (F=F(\bar z)), we have (\partial_{\bar z}F=F'(\bar z)). Therefore
[
2\int_U F,\partial_{\bar z}\phi,dA
==================================

-2\int_U \phi,F'(\bar z),dA
+\frac1i\int_{\partial U}\phi F,dz.
]

With (\phi=e^{-|z|^2/2}), this gives
[
\int_U e^{-|z|^2/2}\big(F'(\bar z)-zF(\bar z)\big),dA
=====================================================

-\int_U e^{-|z|^2/2}F'(\bar z),dA
+\frac1i\int_{\partial U} e^{-|z|^2/2}F(\bar z),dz.
]

So the quadrature identity becomes
[
-\int_U e^{-|z|^2/2}F'(\bar z),dA
+\frac1i\int_{\partial U} e^{-|z|^2/2}F(\bar z),dz
==================================================

\frac c2,F'(0).
]

That is already a boundary-plus-bulk formulation.

---

## 3. Package the bulk and boundary into a single distribution

Define the distribution
[
\mu_U := \frac1i,e^{-|z|^2/2}dz\lfloor_{\partial U}
;-;
e^{-|z|^2/2}\mathbf 1_U,dA,
]
understood by duality on test functions (F(\bar z)) as
[
\langle \mu_U, F'(\bar z)\rangle
:=
\frac1i\int_{\partial U} e^{-|z|^2/2}F(\bar z),dz
-------------------------------------------------

\int_U e^{-|z|^2/2}F'(\bar z),dA.
]

Then the (T=-6) quadrature condition is exactly
[
\langle \mu_U, F'(\bar z)\rangle = \frac c2,F'(0).
]

Equivalently, for all antiholomorphic (F),
[
\frac1i\int_{\partial U} e^{-|z|^2/2}F(\bar z),dz
-------------------------------------------------

# \int_U e^{-|z|^2/2}F'(\bar z),dA

\frac c2,F'(0).
]

So yes: the problem can be reformulated as saying that a certain combination of an **area measure on (U)** and a **boundary measure on (\partial U)** acts like a point mass at the origin on the weighted Landau test class.

---

## 4. A more geometric free-boundary formulation

Now let
[
w:=v,\mathbf 1_U,
]
assuming (v) vanishes outside (U). Then applying (L+6) to (w) as a distribution produces:

* the bulk term ((L+6)v) in (U),
* plus boundary distributions coming from differentiating the indicator (\mathbf 1_U).

This is the usual free-boundary mechanism.

Because (L+6) is second order, the boundary contribution has two parts:

1. a term involving the trace (v|_{\partial U}), coming from second derivatives of (\mathbf 1_U),
2. a term involving the magnetic normal derivative of (v), coming from first derivatives hitting (\mathbf 1_U).

Schematically,
[
(L+6)(v\mathbf 1_U)
===================

\mathbf 1_U (L+6)v

* \mathcal B_1(v),\delta_{\partial U}
* \mathcal B_2(v),\partial_\nu\delta_{\partial U}.
  ]

If we want
[
(L+6)(v\mathbf 1_U)=\mathbf 1_U-c,\delta_0
]
with **no boundary distribution at all**, then we must impose boundary conditions that kill both boundary terms.

Just as for the Laplacian, the natural conditions are
[
v=0 \quad\text{on }\partial U,
\qquad
\partial_\nu^A v =0 \quad\text{on }\partial U,
]
where
[
\partial_\nu^A v := \nu\cdot(\nabla-iA)v.
]

So the clean free-boundary problem is:

[
\begin{cases}
(L+6)v = 1 & \text{in } U\setminus{0},[3pt]
v=0 & \text{on }\partial U,[3pt]
\partial_\nu^A v=0 & \text{on }\partial U,[3pt]
(L+6)v = -c,\delta_0 & \text{at }0.
\end{cases}
]

This is the exact magnetic analogue of the Schwarz-potential formulation.

---

## 5. Where the boundary measure enters if you do not impose both conditions

If you only impose (v=0) on (\partial U), but not (\partial_\nu^A v=0), then the compactly supported extension (w=v\mathbf 1_U) satisfies
[
(L+6)w
======

\mathbf 1_U-c,\delta_0
+
\big(\partial_\nu^A v\big),\delta_{\partial U}
]
up to a convention-dependent constant/sign.

So the residual boundary measure is exactly the magnetic flux across the free boundary.

Likewise, if (v) does not vanish on (\partial U), there is also a (\partial_\nu\delta_{\partial U})-type term.

So the hierarchy is:

* **general extension**: boundary measure + boundary derivative-of-measure;
* **impose (v=0)**: only boundary measure remains;
* **impose (v=0) and (\partial_\nu^A v=0)**: no boundary term remains.

That is the free-boundary interpretation you were asking for.

---

## 6. How the moment identity fits this picture

The integration-by-parts identity you pointed out is precisely what one gets by pairing the residual boundary measure with the Landau-level test functions.

At (T=-6), those test functions are
[
h_F=2e^{-|z|^2/2}(F'(\bar z)-zF(\bar z)).
]

Pairing the boundary defect with (F) gives
[
\frac1i\int_{\partial U} e^{-|z|^2/2}F(\bar z),dz,
]
while the bulk contributes
[
-\int_U e^{-|z|^2/2}F'(\bar z),dA.
]

So your boundary identity is exactly the weak formulation of:

> the boundary defect of the compactly supported extension is a point mass at the origin when tested against the (T=-6) Landau class.

That is the right free-boundary-with-boundary-measure reformulation.

---

## 7. A concise final formulation

For (A=(y,-x)) and (T=-6), the original problem is equivalent in spirit to finding a domain (U) and a function (v) such that
[
v\in C_c(\mathbb C),\qquad
(L+6)v=\mathbf 1_U-c,\delta_0,
]
which can be viewed as the free-boundary problem
[
\begin{cases}
(L+6)v = 1 & \text{in }U\setminus{0},\
v=0,\ \partial_\nu^A v=0 & \text{on }\partial U,\
(L+6)v=-c,\delta_0 & \text{at }0.
\end{cases}
]

If one relaxes the second boundary condition, then one gets instead
[
(L+6)(v\mathbf 1_U)
===================

\mathbf 1_U-c,\delta_0
+
\lambda_{\partial U},
]
where
[
\lambda_{\partial U}\sim (\partial_\nu^A v),\delta_{\partial U}.
]

And when you test this against the (T=-6) Landau class, (\lambda_{\partial U}) produces exactly the boundary integral
[
\frac1i\int_{\partial U} e^{-|z|^2/2}F(\bar z),dz.
]

So yes: the problem really does admit a natural reformulation as a **magnetic free-boundary problem with a boundary measure defect**, and the identity you derived is its weak Landau-level version.

The next useful thing would be to compute that boundary defect **exactly**, with all constants and signs, by expanding (L(v\mathbf 1_U)) in local normal-tangential coordinates.



